\clearpage
\fancyhead[L]{\small CONNOISSANCE DES TEMS.}
\hypertarget{native-542-lower}{}\label{native-542-lower}
\begin{center}\small Göttingische gelehrte Anzeigen. Stück 199. S. 1978..1987. 14. Dezember 1811.\\[3pt]\rule{.72\linewidth}{.3pt}\end{center}

\emph{Connoissance des tems ou des mouvemens célestes, à l'usage des astronomes et des navigateurs pour l'an} 1812, \emph{publiée par le bureau des longitudes. Juillet} 181c. Paris bei \textsc{Courcier}. 415 Seiten in Großoktav.

Die Tafel der Längen und Breiten der wichtigsten Orte der Erde ist in diesem Jahrgang erheblich verbessert worden. Vor allem aus \textsc{Monteiro}s Ephemeriden von Coimbra und aus \textsc{Oltmanns}' Untersuchungen zur Geografie des Neuen Kontinents sind über hundert neue Einträge hinzugekommen, sodass es jetzt insgesamt nahezu 1500 sind. Herr \textsc{Burckhardt} hat auch die Längen von 50 Sternwarten oder anderen durch astronomische Beobachtungen bemerkenswer\SourceLeaf{543}ten Orten in Europa erneut sorgfältig untersucht. Ob die Worte il a comparé et calculé de nouveau \emph{toutes} les observations tant anciennes que modernes ganz wörtlich zu verstehen sind, lassen wir offen. Auf jeden Fall wäre es wünschenswert, dass Herr \textsc{Burckhardt} die Einzelheiten dieser nützlichen Arbeit ausführlicher bekannt gäbe. \textsc{Bouvard}s Beobachtungen an der kaiserlichen Sternwarte im Jahr 1809 machen diesmal beinahe die Hälfte der \emph{Additions} aus. Man weiß, wie wertvoll die Beobachtungen von Sonne, Mond, Planeten und Sternen sind. Wir haben in diesem Jahrgang mit Freude eine Reihe beobachteter Durchgänge des Polarsterns durch den Meridian bemerkt. Schade nur, dass es sich stets um untere Kulminationen handelt. Damit entfällt freilich der vielfältige Nutzen, den zahlreiche Beobachtungen des Polarsterns \emph{in beiden Kulminationen zusammengenommen} bieten können. Für Ceres und Pallas finden sich um die Zeit ihrer Opposition recht viele Beobachtungen. Ebenso finden sich zahlreiche Verfinsterungen von Jupitermonden, aber nur drei unvollständig beobachtete Sternbedeckungen, nämlich der Austritt von $\gamma$ Scorpii am 28. Mai um $12^{\mathrm u}\,3^{\mathrm m}\,53^{\mathrm s}$ mittlerer Zeit, der Austritt von $2\delta$ Tauri am 28. Sept. um $9^{\mathrm u}\,43^{\mathrm m}\,0^{\mathrm s}$ und der Eintritt von $1\delta$ Tauri am 25. Okt. um $18^{\mathrm u}\,17^{\mathrm m}\,15^{\mathrm s}2$. --- Ein einfaches Verfahren, die Positionen des Mondes näherungsweise zu berechnen, von J. C. \textsc{Burckhardt}. Herr B. verwendet dazu die bekannte chaldäische Periode von 18 Jahren oder 223 Lunationen, nach deren Ablauf die Argumente der Mondungleichheiten wieder nahezu ihre früheren Werte annehmen. Er berechnet die Änderungen der Länge und Breite, von denen die beträchtlichsten in fünf Tafeln zusammengestellt sind. Aus einer Ephemeride der Mondbewegung für ein gegebenes Jahr lässt sich so mit sehr wenig Mühe und recht guter Genauigkeit eine ähnliche Ephemeride für ein um eine Periode späteres Jahr berechnen. Wir hätten uns gewünscht, dass Herr \textsc{Burckhardt} auch die Änderung der Horizontalparallaxen auf ähnliche Weise behandelt hätte. --- Über ein neues Verfahren zur Verbesserung der Pendeluhren, von demselben. Um Rost an den Zapfen der Räder zu verhindern, schlägt Herr \textsc{Burckhardt} vor, sie im Feuer zu vergolden und anschließend erneut zu härten. Ob man dabei etwas gewinnt, wird sich am besten zeigen, wenn der Vorschlag tatsächlich ausgeführt wird. --- Tafeln für die Aberration, Nutation und Präzession der 36 Fundamentalsterne von \textsc{Maskelyne}, von demselben. Die Solarnutation ist mit der Aberration zusammengefasst. --- Über die Absenkung des Quecksilbers in Barometerröhren infolge der Kapillarität, von \textsc{Laplace}. Nach \textsc{Laplace}s Theorie bewirkt die Kapillarwirkung des Glases weder die konvexe Oberfläche des Quecksilbers noch die damit verbundene Absenkung in der Röhre. Beides geht auf die Kapillarwirkung einer äußerst dünnen Feuchtigkeitsschicht zurück, die auch bei sorgfältigem Auskochen des Quecksilbers noch an der Glasfläche zurückbleibt. Von der Gestalt der Oberfläche hängt die Beziehung zwischen der Weite der Röhre und der Absenkung des Quecksilbers ab. \textsc{Laplace} hat diese Gestalt durch eine mühsame, auf mechanischer Quadratur beruhende Integration der im Supplement zu Mécanique céleste aufgestellten Grundgleichung bestimmt (unserer Meinung nach hätte es die Rechnung vielleicht noch etwas erleichtert, die Reihenentwicklung etwas weiter zu führen). So hat er für den praktischen Gebrauch eine Tafel berechnet, die hier mitgeteilt wird. Wir hätten uns gewünscht, dass auch die Höhe des konvexen Teils des Quecksilbers, die sich aus dieser Rechnung ergibt, mit angegeben worden wäre. --- Über die (im Septemberheft 1808 der Monatl. Correspondenz) vom \textsc{Gauss} angegebene Lösung einer Aufgabe der sphärischen Astronomie, von \textsc{Delambre}. Aus drei gleichen Höhen bekannter Sterne werden dabei zugleich die Polhöhe des Ortes, der Uhrstand und der Fehler des Instruments bestimmt. Die Behandlung der Aufgabe, die Herr \textsc{Delambre} hier vorlegt, entspricht im \emph{Wesentlichen} der von \textsc{Gauss}, hat aber unserer Meinung nach an Einfachheit und Eleganz verloren. Herr \textsc{Delambre} hat das ganze Beispiel einer Beobachtung, an dem \textsc{Gauss} damals die Rechnung erläuterte, erneut in großer Ausführlichkeit durchgerechnet. Dem von \textsc{Gauss} gefundenen Ergebnis über den Einfluss der Beobachtungsfehler auf die Genauigkeit der Breitenbestimmung \SourceLeaf{544} \[ \mathrm{d}\varphi = 3{,}8077\Delta - 0{,}2884\Delta' + 3{,}5193\Delta'' \] fügt er die Bemerkung bei, Cette dernière formule prouve que dans la pratique la méthode n'aurait qu'une exactitude assez bornée. Dieses unpassende Urteil lässt sich nur dadurch erklären, dass er, kaum begreiflich, übersehen hat, dass $\Delta$, $\Delta'$, $\Delta''$ in Zeitsekunden und $\mathrm{d}\varphi$ in Bogensekunden ausgedrückt sind. Ebenso zeigt der Zusatz, dans une nuit, où l'on pourroit observer trois étoiles à la même hauteur, on verroit très probablement passer au méridien quelque étoile connue qui donneroit la latitude avec moins de peine et plus de précision, et l'heure de la pendule par une simple hauteur, dass Herr \textsc{Delambre} abermals vergessen hat, dass diese Methode für alle Fälle bestimmt ist, in denen man sich bei den absoluten Höhen auf sein Instrument nicht verlassen kann. Was die \emph{Genauigkeit} angeht, ist die Bemerkung völlig unbegründet. --- Es folgen noch einige Ausführungen zu der bekannten, in der Praxis aber so gut wie völlig unbrauchbaren Aufgabe, aus drei Höhen eines einzigen Sterns zugleich seine Deklinationen, seinen Stundenwinkel und die Polhöhe zu bestimmen. Herr \textsc{Delambre} kommt daraufhin noch einmal auf die kleine Abhandlung von \textsc{Gauss} zurück, in der die Beobachtung zweier Höhen zweier bekannter Sterne zur Bestimmung von Zeit und Breite vorgeschlagen worden war und die er im vorhergehenden Band der Connaissance des tems besprochen hatte. In dem, was er darüber und über den Vorzug der Synthese vor der Analyse sagt, sehen wir einen weiteren Beleg für das Urteil, das wir bei der Besprechung jenes Bandes der Connoissance des tems ausgesprochen haben. Herr \textsc{Delambre} hat den Zweck jener Schrift völlig falsch aufgefasst. Dieser in der Schrift selbst ganz klar ausgesprochene Zweck war, eine \emph{Kombination} von Beobachtungen zur Bestimmung der Polhöhe zu empfehlen. Sie ist dazu unter manchen Umständen sehr brauchbar und war, soweit der Verfasser damals wusste und bis heute weiß, dafür in dieser Allgemeinheit noch nicht vorgeschlagen worden. Deshalb nannte er sie eine \emph{neue} Methode. Die \emph{Berechnung} aus solchen Beobachtungen beruht, wie \textsc{Gauss} damals ebenfalls zeigte, auf einem Problem, dessen geometrische Lösung seit \textsc{Tycho}s Zeiten bekannt ist. Nachdem der Verfasser diese mit wenigen Worten vollständig skizziert hatte, verwendete er noch zwei Seiten darauf zu zeigen, dass dieselbe Lösung auch auf rein analytischem Weg kurz und elegant gefunden werden kann. Das tat er nur, weil er glaubte, dass Freunde der Analysis an einer solchen Entwicklung, die sich gerade nicht von selbst ergibt und einiges Geschick erfordert, Gefallen finden könnten. Er hätte sich nicht träumen lassen, dass jemand dies so auslegen könnte, als solle damit jede geometrische Behandlung verdrängt werden, deren Wert hinreichend bekannt und unbestritten ist. Die Bemerkungen von Herrn \textsc{Delambre} über die jeweiligen Vorzüge des analytischen und des geometrischen Verfahrens sind deshalb, obwohl sie meist völlig begründet sind, hier ganz fehl am Platz. --- Über die verschiedenen Verfahren, mit denen Astronomen Sonnenfinsternisse beobachtet haben, von \textsc{Delambre}. Hier wird eine bemerkenswerte Stelle aus \textsc{Appians} Astronomicum Caesareum, gedruckt 1540, angeführt, in der erstmals farbige Gläser für Sonnenbeobachtungen vorgeschlagen werden. Herr \textsc{Delambre} hält es allerdings aus triftigen Gründen für wahrscheinlich, dass \textsc{Appian} seinen Vorschlag selbst nicht umzusetzen versucht hat. --- Darauf folgt ein weitschweifiger, 50 Seiten füllender Auszug aus Gauss' \emph{Theoria motus corporum coelestium}. Herr \textsc{Delambre} hat den größten Teil der darin enthaltenen Formeln herausgeschrieben, meist ohne die Beweise, und die erläuternden Zahlenbeispiele erneut durchgerechnet, wenn auch mit geringerer Genauigkeit als im Werk selbst. Dies hat er mit gelegentlich eingestreuten Anmerkungen hier abdrucken lassen. So empfehlenswert es ist, solche Werke auf diese Weise zu studieren, so wenig scheinen solche Auszüge zum Abdruck geeignet. Wegen des beschränkten Raums können wir hier nur einige der Anmerkungen ansprechen. Die Methode, die Berechnung der geozentrischen Orte der Planeten mithilfe rechtwinkliger Koordinaten unmittelbar auf den Äquator zu beziehen, hat Herr \textsc{Delambre} auf \emph{ein} Beispiel angewandt. Dabei hat er den Vorzug dieses \SourceLeaf{545}Verfahrens gegenüber dem üblichen nicht recht erkennen können. Das widerspricht aber ganz dem Geist dieser Methode, die nur für Fälle bestimmt ist, in denen \emph{viele} geozentrische Orte berechnet werden sollen. Hätte Herr \textsc{Delambre} nach dieser Methode statt eines einzigen Planetenorts ein Dutzend berechnet, so hätte er dasselbe festgestellt wie schon so viele andere Rechner, dass man damit weniger als halb so viel Zeit und Mühe braucht wie mit dem üblichen Verfahren. Die vier Formeln der sphärischen Trigonometrie, die in der Theoria motus so oft verwendet werden, hat Herr \textsc{Delambre} auch selbst gefunden, ihren Vorzug gegenüber dem üblichen Verfahren aber nicht erkannt. Sollen zum Beispiel aus zwei Seiten und dem eingeschlossenen Winkel eines sphärischen Dreiecks alle übrigen Stücke bestimmt werden, so muss man nach den neuen Formeln an sechs verschiedenen Stellen der Sinustafeln insgesamt 12 Logarithmen nachschlagen. Damit erhält man zugleich eine Kontrolle der Rechnung und stets genaue, niemals mehrdeutige Ergebnisse. Beim üblichen, von Herrn \textsc{Delambre} bevorzugten Verfahren muss man hingegen, wenn man ebenfalls eine Kontrolle der Rechnung haben will, insgesamt 13 Logarithmen an 11 verschiedenen Stellen der Tafel nachschlagen. Man erhält dann die dritte Seite aus ihrem Sinus, und zwar, wenn sie nahe am rechten Winkel liegt, ungenau, ja vielleicht sogar mehrdeutig, obwohl sie im Problem selbst eindeutig bestimmt ist. Der Vorzug der neuen Formeln ist somit ganz entschieden und erheblich genug, wenn viele derartige Rechnungen auszuführen sind. Es ist daher verwunderlich, dass Herr \textsc{Delambre} ihn übersehen konnte. Einer Berichtigung bedarf die Aussage, die Herr \textsc{Delambre} auf S.\,357 über die zweite \textsc{Gauss}sche Lösung des Problems macht, aus zwei Abständen eines Planeten von der Sonne, dem eingeschlossenen Winkel und der dazwischenliegenden Zeit die Elemente zu bestimmen. Sie erfordere \emph{sehr lange} Entwicklungen, die der Verfasser nicht angegeben habe, und beruhe teilweise auf seinen eigenen, bisher nicht bekannt gemachten Theorien. Danach sollte man glauben, dass diese Lösung einer der wichtigsten Aufgaben des ganzen Werks so dargestellt ist, dass Leser ohne besondere Übung gar nicht damit zurechtkommen und selbst Kenner keinen ganz vollständigen Einblick in die Lösung gewinnen können. Beides trifft aber nicht zu. Der Rezensent weiß aus vielen Beispielen, dass selbst Anfänger nirgends Schwierigkeiten mit den kleinen, unbedeutenden Entwicklungen hatten, die man stets dem Leser überlassen muss, sofern man kein durch lästige Weitschweifigkeit ungenießbares Buch schreiben will. Mit den eigenen Theorien des Verfassers, deren Entwicklung er sich für eine andere Gelegenheit vorbehalten musste, weil sie hier nicht am Platz war, hängen bei dieser Lösung nur die Rechenvorteile zusammen, die er selbst zur Konstruktion der Hilfstafel genutzt hat und die hier keineswegs wesentlich sind. Für die Hauptgleichung in der großen Aufgabe, die Bahn aus drei geozentrischen Orten zu bestimmen, hatte \textsc{Gauss} die indirekte Lösung als besonders bequem empfohlen, zur Ausführung aber nichts weiter gesagt. Denn zum einen sind die anzuwendenden Kunstgriffe an sich hinreichend bekannt, zum anderen kann die in der Theoria angegebene Lösung des \textsc{Kepler}schen Problems gewissermaßen als Muster dienen. Die Art jedoch, wie Herr \textsc{Delambre} die numerische indirekte Lösung der Hauptgleichung in dem von ihm aufs Ausführlichste erneut durchgerechneten Beispiel angeht, ist mehr als dreimal so lang wie nötig, obwohl er selbst sie noch bequem genug findet. Mit zwei Versuchen kann man nämlich weiter kommen als Herr \textsc{Delambre} mit sieben. Das zeigt, dass es doch nicht unzweckmäßig wäre, wenn der Verfasser gelegentlich an geeigneter Stelle nochmals auf die kleinen Kunstgriffe aufmerksam machte, die bei der indirekten Lösung anzuwenden sind. --- Nach der langen Besprechung der Theoria, bei der sich Herr \textsc{Delambre} auf das erste Buch und den ersten Abschnitt des zweiten beschränkt hat, folgen noch neue Bemerkungen zur Parallaxenrechnung und zu den Formeln der Herren \textsc{Olbers} und \textsc{Littrow}, von \textsc{Delambre}. Darin erklärt er erneut, dass er Lösungen von Aufgaben, die darauf beruhen, die \SourceLeaf{546}Lage von Punkten im Raum auf drei rechtwinklige Koordinaten zu beziehen und die Entwicklungen rein analytisch durchzuführen, stets für weniger einfach hält als solche, die allein sphärische Trigonometrie verwenden. Das ist ein sehr einseitiges Urteil, dem wenige Mathematiker zustimmen werden. Eine Formel, in der die Längenparallaxe durch eine Reihe nach den Sinus der Vielfachen des Abstands des Mondes vom Nonagesimus dargestellt wird, hatte \textsc{Olbers} Herrn \textsc{Rohde} zugeschrieben. Herr \textsc{Delambre} beansprucht sie hier als seine eigene Erfindung, die er schon in der Connoissance des tems für 1793 bekannt gemacht habe. Eigentlich gebührt sie jedoch \textsc{Lagrange}, der sie bereits in etwas allgemeinerer Form in den Nouveaux Mémoires de l'académie de Berlin, 1776, S. 231, bekannt gemacht hatte. --- Ein Verfahren, mit dem sich eine Uhr zur Anzeige von Sternzeit und mittlerer Zeit verwenden lässt, von \textsc{Burckhardt}, beruht auf dem angenäherten Verhältnis beider Zeiten, 51.79 zu 49.82, das in einem ganzen Jahr nur um 4 Sekunden abweicht. Von \emph{demselben}, über den zweiten Kometen von 1737, nach Beobachtungen in China, die in der Monatl. Corresp. bekannt gemacht worden sind. --- Die meteorologischen Beobachtungen an der Pariser kaiserlichen Sternwarte im Jahr 1808 und das Verzeichnis der Mitglieder des Längenbüros bilden wie gewöhnlich den Schluss des Bandes.
